\section{Background and Motivation}
\label{sec:background}

\subsection{Heuristic JML inference and single-pass propagation}

The inferrer of~\cite{edmonds2026inference} is a heuristic JML producer for Java in the tradition of contract-by-design~\cite{meyer1992design,jml,leavens2006design,chalin2010jml}: a library of pattern-matchers over the JavaParser~\cite{javaparser} abstract syntax tree (AST) that emits a \texttt{MethodSpecification} per method declaration. Twelve specialist analysers (precondition, postcondition, loop-invariant, assignable, nullability, overflow, etc.) cooperate via a \texttt{SpecificationCache} indexed by simple-class method signature. The semantic target is OpenJML~\cite{openjml} extended static checking, which traces its lineage to ESC/Java2~\cite{flanagan2002extended} and the broader pluggable-types programme~\cite{papi2008practical}.

Cross-method propagation is handled by an
\texttt{InterproceduralAnalyzer} component invoked during isolated
inference. When method \texttt{m} contains a call \texttt{n(args)},
this analyser looks up \texttt{n}'s cached specification, substitutes
formal parameters with actual arguments, and appends the substituted
preconditions to \texttt{m}'s precondition set. The pass is
single-pass: each method is visited once in topological order over the
call graph; the callee's spec is read once and added at the call site
without further transformation. We refer to this as the \emph{isolated}
strategy throughout the paper, since each method's analysis treats
callees as black boxes whose contracts are pulled in atomically.

\subsection{Where the single-pass strategy is blind}

The single-pass propagation is the right strategy for the simple case:
a precondition that the callee requires \emph{unconditionally} must
hold at the call site, and conjoining it into the caller's set is
sound. But the strategy is structurally blind to two classes of
information.

\textbf{Branch-conditional obligations.} If the caller \texttt{m} has a
body fragment
\begin{lstlisting}[language=Java,frame=none]
if (g) {
    n(args);  // n requires p
}
\end{lstlisting}
then \texttt{m}'s precondition is not \texttt{p[args/params]} but
rather \texttt{g ==> p[args/params]} --- the obligation must hold only
on the path where the call is reached. Single-pass propagation either
appends \texttt{p[args/params]} unconditionally (overconstraining
\texttt{m} on paths where the call does not fire) or skips the
substitution to be safe (missing the obligation entirely on the
guarded path). Neither is correct.

\textbf{Polymorphic-dispatch obligations.} If \texttt{m} calls
\texttt{x.n()} where \texttt{x}'s declared type has multiple
implementations of \texttt{n}, runtime dispatch may pick any
candidate, each with its own inferred precondition; single-pass
propagation resolves the call to one nominal callee and sees only
one. The compositional pass emits the disjunction
\texttt{(p\_1 || \dots\ || p\_k)} over the candidates' substituted
preconditions. We treat this as a deliberate \emph{weakening} of the
caller obligation, not a sound safety guarantee: since the
heuristically inferred per-override preconditions need not respect
behavioural subtyping, the genuinely sound obligation (the
conjunction, or a type-guarded case split) is frequently spuriously
strong, whereas the weakest obligation satisfiable by some candidate
is the more useful signal. Either way it is a shape a single nominal
resolution cannot produce.

\subsection{A motivating example}
\label{sec:motivating}

Consider a caller \texttt{m} that calls a callee \texttt{n} only on one branch:

\begin{lstlisting}
//@ requires divisor != 0;
static int safeDivide(int a, int divisor) { return a / divisor; }

static int scale(int value, int factor) {
    if (factor != 0) {
        return safeDivide(value, factor);
    }
    return value;
}
\end{lstlisting}

The single-pass strategy faces a dilemma at the call site
\texttt{safeDivide(value, factor)}. Substituting yields the clause
\texttt{factor != 0}; conjoining it into \texttt{scale}'s precondition would
forbid the very call \texttt{scale(value, 0)} that the method explicitly
handles, while skipping it loses the obligation on the path where the call
fires. Neither is correct. The compositional pass walks \texttt{scale}'s body
and lifts the substituted clause through the enclosing guard, producing the
correct \emph{branch-conditional} clause:

\begin{lstlisting}[language={}]
//@ requires factor != 0 ==> factor != 0;   // tautology here
\end{lstlisting}

The example is deliberately simple so the lifting is visible; in the corpus
the guard and the callee's requirement differ (\textit{e.g.}, a call to a
method requiring a non-empty list guarded by a size check), and the lifted
clause \texttt{size > 0 ==> list.size() > 0} carries genuine, non-tautological
information the single pass cannot express. The polymorphic-dispatch
disjunction (above) is the second such shape. These two shapes are the
structural subjects of Section~\ref{sec:empirical}.

\subsection{What a second-pass compositional analysis would add}

The classic weakest-precondition calculus of
Dijkstra~\cite{dijkstra_guarded_1975} addresses both shapes
structurally. Walking a method's body in WP fashion accumulates the
guard context at each program point, so a call site reached only on
\texttt{g} contributes \texttt{g ==> p}; a polymorphic call site
contributes the disjunction over candidates by construction.

Implementing a full WP transformer for arbitrary Java would require
a satisfiability-modulo-theories (SMT) decision procedure and a precise operational semantics for every
language construct, which is the domain of full verification tools
(KeY, OpenJML's SC mode) rather than heuristic inferrers. A
\emph{scaffold} that walks the body in a WP-flavoured order, lifts
each call site's substituted callee specification through enclosing
branch guards, and recognises polymorphic candidates is, however,
implementable as an approximately 450-line pass over the AST. That is what this
paper does.

\subsection{What we measure, and what we do not}

We do not measure whether the second-pass specifications are
\emph{semantically correct} against the code's intent --- that is the
inferrer's responsibility, not the compositional pass's, and it is
the question of~\cite{edmonds2026inference}. The compositional pass
is a structural transformation: given the cached callee specs
(whatever their semantic quality), it produces a new spec for the
caller. We measure whether that new spec contains clauses the single
pass did not.

The downstream impact on LLM-based test generation (whether the
richer specs measurably improve mutation-killing rate) is the subject
of a planned follow-up. The present article is scoped to the
structural claim --- does the second pass add anything to a real-world
single-pass baseline --- because the answer matters even if the
downstream impact is small (it tells us the architectural distinction
is real) and because the answer is empirically tractable in a way the
downstream comparison is not.
